\section{Discussion and Future Work}
\label{sec:disc}

\subsection{Answer to Research Questions}~\label{sec:disc:rq}

We answer the research questions as follows:

\begin{itemize}
  \item \textbf{RQ1:} Our parallel congruence closure algorithm achieves a sublinear,
  but meaningful speedup on random benchmarks, a synthetic benchmark suite
  designed to stress test our algorithm, and a set of equality saturation
  benchmarks~\cite{eqsat}. We observe meaningful speedups up to 32 cores, but we experienc
  a pleateau on 32 and then 64 cores.
  \item \textbf{RQ2:} Our algorithm tends to perform better on benchmarks with a large
  number of merges per round. Having more rounds of congruence closure does not necessarily
  lead to worse performance, but it does not lead to better performance either as we only 
  parallelize within rounds.
  \item \textbf{RQ3:} Our algorithm provides a meaningful speedup over a
  sequential implementation on the egg benchmarks. Future work should evaluate
  our algorithm on benchmarks from SMT solvers and  more challenging equality
  saturation benchmarks such as \texttt{egg-cc}~\cite{eggcc}.
\end{itemize}



\subsection{Limitations of Evaluation}



\subsection{Limitations of Algorithm}

- reliance on bulk-synchronous parallelism; a truly concurrent algorithm could
be more efficient and have better speedups, but is more difficult to design and
implement

- non backtrackability

- does not support proof generation, which is important for SMT solvers

- does not support incremental updates, which is important for SMT solvers and
equality saturation engines

\subsection{Potential Future Work}

- truly concurrent algorithm that doesn't rely on bulk-synchronous parallelism